% GENERATED FILE. Edit canonical lessons, then run npm run generate:books.
% canonical-source-hash: fnv1a64:206d14b5c6646582
% canonical-lessons: GU-C347-kaydo, GU-C347-ullanghan, GU-C347-palan, GU-C347-lagu, GU-C347-uthapvu

\chapter{The Law}
\label{ch:gu-the-law}
\hlchaptermodality{347}

\begin{chapteropening}
\textbf{By the end of this chapter:} \emph{I can say what the law allows and what it forbids.}

Complete the last lesson of chapter 347: I can say what the law allows and what it forbids.
\end{chapteropening}

\section[\texorpdfstring{\gu{કાયદો} (kāydo)}{કાયદો (kāydo)}]{\gu{કાયદો} (kāydo) — a law}

\label{lesson:GU-C347-kaydo}



\begin{quote}
Before the new one: say the Gujarati for different, then the Gujarati for to compare.
\end{quote}

\subsection*{You'll want to know: \gu{કાયદો}}
\textbf{\gu{કાયદો}} — \emph{kāydo} — \textquotedblleft{}a law\textquotedblright{}.

\begin{quote}
\gu{નવો} \gu{કાયદો} \gu{ગરીબ} \gu{માણસો} \gu{માટે} \gu{સારો} \gu{છે}.
\end{quote}

\begin{quote}
\emph{navo kāydo garīb māṇso māṭe sāro che.}
\end{quote}

\begin{quote}
The new law is good for poor people.
\end{quote}

A masculine noun.

\begin{cousinweb}
\textbf{\gu{કાયદો}} (\emph{kāydo}) is usually traced to Arabic \emph{qāʿida}, \textquotedblleft{}a base, a rule\textquotedblright{}, through Persian.
\end{cousinweb}

\begin{grammarlens}[title={What you've built}]
One more word for giving reasons and opinions.
\end{grammarlens}

\subsection*{Your turn}
\begin{itemize}
\raggedright
  \item \emph{Say it:} \emph{kāydo}
  \item \emph{Say it:} \emph{kāydo}, once more
  \item \emph{Say it:} \emph{navo kāydo garīb māṇso māṭe sāro che.} — The new law is good for poor people.
  \item \emph{Recall:} say the Gujarati for different, then the Gujarati for to compare, then say \emph{kāydo} again
\end{itemize}

\begin{tcolorbox}[breakable,colback=teal!4,colframe=teal!35!black,title={Before you move on}]
What does \gu{કાયદો} mean? (\textquotedblleft{}A law\textquotedblright{}.) Say it once more.
\end{tcolorbox}
\section[\texorpdfstring{\gu{ઉલ્લંઘન} (ullanghan)}{ઉલ્લંઘન (ullanghan)}]{\gu{ઉલ્લંઘન} (ullanghan) — a violation}

\label{lesson:GU-C347-ullanghan}



\begin{quote}
Before the new one: say the Gujarati for to compare, then the Gujarati for a law.
\end{quote}

\subsection*{You'll want to know: \gu{ઉલ્લંઘન}}
\textbf{\gu{ઉલ્લંઘન}} — \emph{ullanghan} — \textquotedblleft{}a violation\textquotedblright{}.

\begin{quote}
\gu{લાલ} \gu{સિગ્નલ} \gu{ઓળંગવું} \gu{એ} \gu{કાયદાનું} \gu{ઉલ્લંઘન} \gu{છે}.
\end{quote}

\begin{quote}
\emph{lāl signal oḷangvũ e kāydānũ ullanghan che.}
\end{quote}

\begin{quote}
Going through a red light is a breach of the law.
\end{quote}

A neuter noun.

\begin{cousinweb}
\textbf{\gu{ઉલ્લંઘન}} (\emph{ullanghan}) is a tatsama, borrowed from Sanskrit \emph{ullaṅghana}, \textquotedblleft{}leaping over, passing beyond\textquotedblright{} (Monier-Williams): stepping over a line one should not cross.
\end{cousinweb}

\begin{grammarlens}[title={What you've built}]
One more word for giving reasons and opinions.
\end{grammarlens}

\subsection*{Your turn}
\begin{itemize}
\raggedright
  \item \emph{Say it:} \emph{ullanghan}
  \item \emph{Say it:} \emph{ullanghan}, once more
  \item \emph{Say it:} \emph{lāl signal oḷangvũ e kāydānũ ullanghan che.} — Going through a red light is a breach of the law.
  \item \emph{Recall:} say the Gujarati for to compare, then the Gujarati for a law, then say \emph{ullanghan} again
\end{itemize}

\begin{tcolorbox}[breakable,colback=teal!4,colframe=teal!35!black,title={Before you move on}]
What does \gu{ઉલ્લંઘન} mean? (\textquotedblleft{}A violation\textquotedblright{}.) Say it once more.
\end{tcolorbox}
\section[\texorpdfstring{\gu{પાલન} (pālan)}{પાલન (pālan)}]{\gu{પાલન} (pālan) — observance (keeping a rule)}

\label{lesson:GU-C347-palan}



\begin{quote}
Before the new one: say the Gujarati for a law, then the Gujarati for a violation.
\end{quote}

\subsection*{You'll want to know: \gu{પાલન}}
\textbf{\gu{પાલન}} — \emph{pālan} — \textquotedblleft{}observance (keeping a rule)\textquotedblright{}.

\begin{quote}
\gu{દરેકે} \gu{નિયમોનું} \gu{પાલન} \gu{કરવું} \gu{જોઈએ}.
\end{quote}

\begin{quote}
\emph{dareke niyamonũ pālan karvũ joīe.}
\end{quote}

\begin{quote}
Everyone should follow the rules.
\end{quote}

A neuter noun.

\begin{cousinweb}
\textbf{\gu{પાલન}} (\emph{pālan}) is a tatsama, borrowed from Sanskrit \emph{pālana}, \textquotedblleft{}guarding, keeping\textquotedblright{} (Monier-Williams).
\end{cousinweb}

\begin{grammarlens}[title={What you've built}]
One more word for giving reasons and opinions.
\end{grammarlens}

\subsection*{Your turn}
\begin{itemize}
\raggedright
  \item \emph{Say it:} \emph{pālan}
  \item \emph{Say it:} \emph{pālan}, once more
  \item \emph{Say it:} \emph{dareke niyamonũ pālan karvũ joīe.} — Everyone should follow the rules.
  \item \emph{Recall:} say the Gujarati for a law, then the Gujarati for a violation, then say \emph{pālan} again
\end{itemize}

\begin{tcolorbox}[breakable,colback=teal!4,colframe=teal!35!black,title={Before you move on}]
What does \gu{પાલન} mean? (\textquotedblleft{}Observance (keeping a rule)\textquotedblright{}.) Say it once more.
\end{tcolorbox}
\section[\texorpdfstring{\gu{લાગુ} (lāgu)}{લાગુ (lāgu)}]{\gu{લાગુ} (lāgu) — in force; applicable}

\label{lesson:GU-C347-lagu}



\begin{quote}
Before the new one: say the Gujarati for a violation, then the Gujarati for observance.
\end{quote}

\subsection*{You'll want to know: \gu{લાગુ}}
\textbf{\gu{લાગુ}} — \emph{lāgu} — \textquotedblleft{}in force; applicable\textquotedblright{}.

\begin{quote}
\gu{નવો} \gu{કાયદો} \gu{આવતા} \gu{મહિનાથી} \gu{લાગુ} \gu{થશે}.
\end{quote}

\begin{quote}
\emph{navo kāydo āvtā mahināthī lāgu thaśe.}
\end{quote}

\begin{quote}
The new law comes into force next month.
\end{quote}

\begin{cousinweb}
\textbf{\gu{લાગુ}} (\emph{lāgu}) CDIAL lists Gujarati \emph{lāgu}, \textquotedblleft{}near\textquotedblright{}, under Sanskrit \emph{lagna}, \textquotedblleft{}sticking to, fixed on\textquotedblright{} (CDIAL 10893); \textquotedblleft{}in force, applicable\textquotedblright{} is its sense in law and in offices.
\end{cousinweb}

\begin{grammarlens}[title={What you've built}]
One more word for giving reasons and opinions.
\end{grammarlens}

\subsection*{Your turn}
\begin{itemize}
\raggedright
  \item \emph{Say it:} \emph{lāgu}
  \item \emph{Say it:} \emph{lāgu}, once more
  \item \emph{Say it:} \emph{navo kāydo āvtā mahināthī lāgu thaśe.} — The new law comes into force next month.
  \item \emph{Recall:} say the Gujarati for a violation, then the Gujarati for observance, then say \emph{lāgu} again
\end{itemize}

\begin{tcolorbox}[breakable,colback=teal!4,colframe=teal!35!black,title={Before you move on}]
What does \gu{લાગુ} mean? (\textquotedblleft{}In force; applicable\textquotedblright{}.) Say it once more.
\end{tcolorbox}
\section[\texorpdfstring{\gu{ઉથાપવું} (uthāpvũ)}{ઉથાપવું (uthāpvũ)}]{\gu{ઉથાપવું} (uthāpvũ) — to defy (an order or a word)}

\label{lesson:GU-C347-uthapvu}



\begin{quote}
Before the new one: say the Gujarati for observance, then the Gujarati for in force.
\end{quote}

\subsection*{You'll want to know: \gu{ઉથાપવું}}
\textbf{\gu{ઉથાપવું}} — \emph{uthāpvũ} — \textquotedblleft{}to defy (an order or a word)\textquotedblright{}.

\begin{quote}
\gu{તેણે} \gu{પિતાની} \gu{વાત} \gu{ઉથાપી}.
\end{quote}

\begin{quote}
\emph{teṇe pitānī vāt uthāpī.}
\end{quote}

\begin{quote}
He went against his father's word.
\end{quote}

Past, with \gu{તેણે} (he or she …): \gu{તેણે} \gu{નિયમ} \gu{ઉથાપ્યો}, \emph{teṇe niyam uthāpyo}.

\begin{cousinweb}
\textbf{\gu{ઉથાપવું}} (\emph{uthāpvũ}) is inherited from a form CDIAL reconstructs as \emph{ut-sthāpyatē}, \textquotedblleft{}is raised\textquotedblright{}; CDIAL glosses the Gujarati verb \textquotedblleft{}to take down, to remove, to transgress\textquotedblright{} (CDIAL 1904).
\end{cousinweb}

\begin{grammarlens}[title={What you've built}]
One more word for giving reasons and opinions.
\end{grammarlens}

\subsection*{Your turn}
\begin{itemize}
\raggedright
  \item \emph{Say it:} \emph{uthāpvũ}
  \item \emph{Say it:} \emph{uthāpvũ}, once more
  \item \emph{Say it:} \emph{teṇe pitānī vāt uthāpī.} — He went against his father's word.
  \item \emph{Recall:} say the Gujarati for observance, then the Gujarati for in force, then say \emph{uthāpvũ} again
\end{itemize}

\begin{tcolorbox}[breakable,colback=teal!4,colframe=teal!35!black,title={Before you move on}]
What does \gu{ઉથાપવું} mean? (\textquotedblleft{}To defy (an order or a word)\textquotedblright{}.) Say it once more.
\end{tcolorbox}
